\section{Questions, and what was not checked}\label{sec:open}

\subsection{Questions}

\begin{enumerate}[label=\arabic*.,ref=\arabic*]
\item\label{q:global} \emph{The global statement.} Show that for every prime $p\equiv1,121,169,289,361,529\pmod{840}$ some shell is occupied. This is the conjecture itself, restated by Theorem~\ref{thm:shell}, and the archive names it as the remaining obligation (p.~4). By Proposition~\ref{prop:groupbarrier}, a proof must produce, for some $R$, a shell whose group $H_a$ contains $-1$ or $-4$; when $R$ is prime this means a prime factor of $(p+R)/4$ that is a non-residue modulo $R$. That condition alone does not suffice, and Proposition~\ref{prop:groupconverse} says when it does.
\item\label{q:typeII} \emph{Type II universality} (Section~\ref{ssec:locks}). Does every prime $p\equiv1\pmod{24}$ have a solution with $p$ dividing exactly two denominators? A positive answer proves the conjecture. Below $10^7$ every such prime has one in a shell $R\le107$, and only $p_*=8{,}803{,}369$ needs $107$; a reading pass found the same below $10^8$.
\item\label{q:qrsurvivors} \emph{The quadratic-residue survivors} (Table~\ref{tab:survivors}). Are there infinitely many primes that satisfy every condition of Section~\ref{ssec:qrconditions}? They have density zero (Proposition~\ref{prop:densityzero}), and $23$ of them lie below $10^7$. A negative answer would prove the conjecture at all but finitely many primes.
\item\label{q:550} \emph{The $550$ supported classes} (Theorem~\ref{thm:termcore}). Is each of the $550$ non-square classes left open modulo $840\cdot11\cdot19\cdot23$ covered by finitely many single identities at some finite level? If one is not, no identity sieve reduces the problem to the square core.
\item\label{q:census} \emph{The census survivors.} $2.394\%$ of the final hard domain survives the carrier chain (§\ref{ssec:census}), and the square rows among them can never be forced by its carriers (Corollary~\ref{cor:censuslimit}). Can the non-square survivor rows all be forced by carriers of the same kind, if more moduli are added? Which carriers, using the factorisation of $(p+R)/4$ and not only residues of $p$, force the square rows?
\item\label{q:deficit} \emph{The strict-record deficit-box conjecture} (§\ref{ssec:records}; archive §19.2). It holds at the two strict records with non-empty deficit below $8{,}803{,}369$, and strict recordness cannot be dropped from it ($p=806{,}521$). A further test needs a strict record above $8{,}803{,}369$, and there is none below $3\cdot10^9$ (Proposition~\ref{prop:records}).
\item\label{q:growth} \emph{Growth of the least class.} Is $h$ unbounded on the primes? Along the strict records up to $3\cdot10^9$ it takes the values $1,2,3,6,8,15,27$; below $3\cdot10^9$ the next largest values are $21$, $20$ and $19$, at nine primes; and every record after $73$ lies in a hard class (Corollary~\ref{cor:hle2}). The workbench states no conjecture on the growth of $h(p)$, and none is stated here.
\item\label{q:shear} \emph{The density of the shear paths.} What is the exact relative density, among the primes $p\equiv2\pmod3$, of the primes with a path of two shear edges? It is at least $1/24$, from the class $131\bmod180$, and at most $0.335$ (Remark~\ref{rem:sheardensity}); below $10^7$ the proportion is $5.24\%$.
\item\label{q:genera} \emph{Genera of the monodromy covers.} What are the genera of the components of the $(3,4,\infty)$ covers at every level? Theorem~\ref{thm:orbits} gives the orbits, and the genera follow from the cycle counts of $A_1$, $A_2$ and $A_\infty$ on each orbit (§\ref{ssec:es-s6}).
\item\label{q:brieskorn} \emph{Brieskorn spheres $\Sigma(2,3,q)$.} For which $q$ does the shadow representation of the associated mock vector coincide, after conjugation and a twist, with a minimal-model representation, as for $q=7$ (Theorem~\ref{thm:m27}), or with a Galois conjugate of one, as for $q=5$ (Proposition~\ref{prop:fibonacci})?
\item\label{q:modularobject} \emph{A modular object for the shell counts.} Is there a modular, mock or quantum modular object whose coefficients are the shell counts $2|E_a|+|M_a|$, or a typed map from the pair-count series of Proposition~\ref{prop:pairzeta} to the mock vector of Section~\ref{ssec:mock}?
\item\label{q:thin} \emph{A typed map to the Descartes geometry.} Is there a typed map from shell occupancy, or from the notes' Descartes quadruples, to the Apollonian group? The companion note \cite{ApollonianNote} determines how the rotations and boosts of Proposition~\ref{prop:flowboost} meet that group; the map from the shell data is what is missing.
\item\label{q:pi1} \emph{The least $\Pi_1$ machine.} What is the least number of states of a Turing machine that halts if and only if the conjecture fails?
\end{enumerate}

\subsection{What this paper did not check}

\begin{itemize}
\item \emph{The archive.} Apart from the parts marked proved or reproduced in §\ref{ssec:archive}, its 619 pages were mapped, not audited. In particular, the structural constructions (§§2, 4, 6--8, 9 after its opening, 21--23, the operator part of §24, the structural material of §24.4, and §25) were not audited, and the literature audits (§§11--14 and 17) were not compared with the papers they audit.
\item \emph{Lean.} The source archive contains 67 files with the extension \texttt{.lean} (seven of them \texttt{.axioms.lean}); they were not rebuilt. The archive's own account of what each Lean module checks is quoted where it matters (§\ref{ssec:census}).
\item \emph{Section~\ref{sec:notes}.} Several statements there rest on the reading passes and are marked so:
\begin{itemize}
\item the lattice and modular facts of Section~\ref{ssec:structfacts}, apart from the involution count;
\item the comparison of the notes' $q$-expansions with Zagier's vector $M_7$;
\item the Busy-Beaver residue test, and the Busy-Beaver values that the archive quotes from the literature;
\item the Cayley--Dickson identities;
\item the counts below $10^8$.
\end{itemize}
Salez's completeness statement was not proved here. Drafts that the author's latest record supersedes were not treated.
\item \emph{The census.} The enumerations of §\ref{ssec:census} were not re-run. Only the arithmetic consistency of the reported counts was checked, together with the base rows and the carrier sets used in Corollary~\ref{cor:censuslimit}.
\item \emph{The supplement.} The general proofs of the supplement's Sections 1--7 were not re-derived, except that of the no-two-shears theorem of Section 5 (Proposition~\ref{prop:twoshears}); their finite content was reproduced in part (§\ref{ssec:supplement}). Section 8, the auxiliary torus, was not checked.
\item \emph{The 124-page reader.} Only Proposition~\ref{prop:quartic} was checked.
\item \emph{The continuation.} Items 1, 5, 9--13, 15--20, 24, 25, 27 and 28 were located, not checked. So were item 4's enumeration, the deduction in item 23, and ES-04 to ES-08 beyond the ES-08 witness.
\item \emph{Novelty.} No literature search for novelty was made here, except where a status line says so. The workbench claims no novelty for the classical coordinates and identities it uses. Several results were found by referee passes (Propositions~\ref{prop:sheargraph}, \ref{prop:groupbarrier}, \ref{prop:groupconverse}, \ref{prop:fixeddivisor} and~\ref{prop:nosquare}, and Corollaries~\ref{cor:hle2} and~\ref{cor:censuslimit}); some of them are forms of statements of the archive or the notes, as their status lines say. The results marked new in version~4 are not in the workbench, the notes or the earlier versions of this paper, as far as these were read; whether they are in the literature was not checked, and none is claimed as new mathematics beyond that.
\end{itemize}
